\begin{casebox}
\textbf{知识点小案例：stable filter 为什么要 count、scan、write 三步。}
特例是输入 \code{[a,b,c,d]}，保留 \code{b,d}。串行 \code{push\_back} 很自然；并行时若多个线程同时 push，就会共享写输出尾指针，还可能打乱顺序。三阶段做法先让每个 block 统计保留数量，再 exclusive scan 得到每个 block 的输出起点，最后各 block 写入自己独占区间。由此推出规律：scan 的作用不是数学炫技，而是把不确定的共享追加变成确定的写区间。工程上任何并行压缩、分区和稳定输出都要先定义输出位置，再谈并行写。
\end{casebox}
